% ForestFirePred_presentation.tex
% Beamer presentation template based on ForestFirePred.ipynb analyses

\documentclass[aspectratio=169,11pt]{beamer}

% Design (built-in, higher visual impact)
\usetheme{Warsaw}
\setbeamertemplate{navigation symbols}{}
\setbeamertemplate{footline}[frame number]
\setbeamerfont{title}{series=\bfseries,size=\LARGE}
\setbeamerfont{subtitle}{series=\mdseries,size=\large}
\setbeamerfont{frametitle}{series=\bfseries}

\usepackage[T1]{fontenc}
\usepackage[utf8]{inputenc}
\usepackage[french]{babel}
\usepackage{graphicx}
\usepackage{booktabs}
\usepackage{amsmath}
\usepackage{mathtools}
\usepackage{hyperref}
\usepackage{ifthen}

% Ignore tiny box overflows (theme-dependent, visually negligible)
\AtBeginDocument{\vfuzz=2pt\relax\hfuzz=2pt\relax}

% Color palette (simple + punchy)
\definecolor{FFDark}{HTML}{0B1F3A}
\definecolor{FFAccent}{HTML}{4169E1}
\setbeamercolor{structure}{fg=FFAccent}
\setbeamercolor{frametitle}{bg=FFDark,fg=white}
\setbeamercolor{title}{fg=white}
\setbeamercolor{subtitle}{fg=white}
\setbeamercolor{title page}{bg=FFDark,fg=white}
\setbeamercolor{section in toc}{fg=FFDark}

% Math notation (cleaner, course-like)
\newcommand{\Area}{\ensuremath{\mathrm{area}}}
\newcommand{\AreaCol}{\texttt{area}}
\newcommand{\LogArea}{\ensuremath{\log\!\bigl(1+\Area\bigr)}}
\DeclareRobustCommand{\RMSE}{\ensuremath{\operatorname{RMSE}}}
\DeclareRobustCommand{\MAE}{\ensuremath{\operatorname{MAE}}}

% Keep images inside the slide box (avoid overflow / shifting layouts)
\setkeys{Gin}{keepaspectratio}
\newcommand{\FFigure}[2][]{%
  \centering\includegraphics[width=\linewidth,height=0.58\textheight,#1]{#2}%
}
\newcommand{\FFigureTwo}[3][]{%
  \begin{columns}[T,onlytextwidth]
    \column{0.49\textwidth}
      \centering\includegraphics[width=\linewidth,height=0.46\textheight,#1]{#2}
    \column{0.49\textwidth}
      \centering\includegraphics[width=\linewidth,height=0.46\textheight,#1]{#3}
  \end{columns}
}

\newcommand{\GroupName}{Groupe ForestFire}

\title[Forest Fires]{TP / Mini-projet — Pr\'ediction des feux de for\^et}
\subtitle{R\'egression de la surface br\^ul\'ee et variante classification}
\author[\GroupName]{\texorpdfstring{\GroupName\\[0.2em]
Alban Komla HENOVI\\
Asta Walo DIOP\\
Arthur BOMEMO Wal\'e}{\GroupName}}
\institute{Cours Python — Semestre 1}
\date{\today}

% PDF metadata (avoid line breaks in PDF strings)
\hypersetup{pdfauthor={\GroupName},pdftitle={Forest Fire Prediction}}
\pdfstringdefDisableCommands{\def\\{ } }

\begin{document}

\begin{frame}[plain]
  \titlepage
\end{frame}

\begin{frame}{Plan}
  \begin{enumerate}
    \item Contexte scientifique (FWI)
    \item Donn\'ees et EDA (distribution, corr\'elations)
    \item M\'ethodologie (pr\'e-traitements, split, validation)
    \item R\'egression : r\'esultats et discussion
    \item Classification binaire : r\'esultats et discussion
    \item Bonus et conclusion
  \end{enumerate}
\end{frame}

\section{Contexte}

\begin{frame}{Contexte : indices FWI}
  Les indices \textbf{FWI (Fire Weather Index)} r\'esument l'\og \`etat du combustible\fg{} et le danger de propagation :
  \begin{itemize}
    \item \texttt{FFMC} : humidit\'e des combustibles fins (court terme)
    \item \texttt{DMC} : humidit\'e couche organique (moyen terme)
    \item \texttt{DC} : s\'echeresse profonde (long terme)
    \item \texttt{ISI} : potentiel de propagation initiale (d\'epend du vent)
  \end{itemize}
  Objectif : pr\'edire la surface br\^ul\'ee $\Area$ (colonne \AreaCol, en hectares) \`a partir de la m\'et\'eo et des indices FWI.
\end{frame}

\section{Donn\'ees}

\begin{frame}{Donn\'ees : Forest Fires (\texttt{forestfires.csv})}
  \begin{itemize}
    \item Cible : \AreaCol (surface br\^ul\'ee $\Area$, en ha)
    \item Variables : coordonn\'ees (\texttt{X,Y}), calendrier (\texttt{month,day}), FWI (\texttt{FFMC,DMC,DC,ISI}), m\'et\'eo (\texttt{temp,RH,wind,rain})
    \item Particularit\'e : beaucoup de z\'eros sur \AreaCol + quelques valeurs extr\^emes
  \end{itemize}

  \vspace{0.6em}
  \textit{Remarque :} \texttt{month} et \texttt{day} sont encod\'es en one-hot.
\end{frame}

\begin{frame}{EDA : distribution de \AreaCol et effet de \(\LogArea\)}
  \begin{itemize}
    \item \AreaCol est tr\`es asym\'etrique (beaucoup de z\'eros)
    \item La transformation \(\LogArea\) rend la cible plus exploitable pour la mod\'elisation
  \end{itemize}

  \vspace{0.4em}
  \IfFileExists{figures/dist_area_log.png}{
    \FFigure{figures/dist_area_log.png}
  }{\textit{(Figure non disponible : ex\'ecuter la cellule Export Beamer)} }
\end{frame}

\begin{frame}{Transformation de la cible}
  Comme la distribution est tr\`es asym\'etrique, on utilise :
  \[
    y = \LogArea
  \]
  \begin{itemize}
    \item Stabilise l'apprentissage (moins domin\'e par les extr\^emes)
    \item G\`ere naturellement \(\Area=0\)
    \item Les m\'etriques (\RMSE/\MAE) sont rapport\'ees en espace log, coh\'erent avec l'entra\^inement
  \end{itemize}
\end{frame}

\begin{frame}{EDA : relations simples (temp\'erature / vent)}
  \begin{itemize}
    \item Pearson : relation lin\'eaire (sensible aux outliers)
    \item Spearman : relation monotone (plus robuste)
  \end{itemize}

  \vspace{0.3em}
  \textit{Observation :} corr\'elations souvent faibles (ph\'enom\`ene multi-factoriel et bruit\'e).

  \vspace{0.4em}
  \textbf{Illustration :} nuages de points export\'es depuis le notebook.
  
  \vspace{0.3em}
  \IfFileExists{figures/scatter_temp_area.png}{
    \IfFileExists{figures/scatter_wind_area.png}{
      \FFigureTwo{figures/scatter_temp_area.png}{figures/scatter_wind_area.png}
    }{\FFigure{figures/scatter_temp_area.png}}
  }{\textit{(Figures non disponibles : ex\'ecuter la cellule Export Beamer)} }
\end{frame}

\section{M\'ethodologie}

\begin{frame}{M\'ethodologie : pr\'e-traitements et pipeline}
  \begin{itemize}
    \item Variables num\'eriques : \textbf{standardisation} (important pour SVR/KNN)
    \item Variables cat\'egorielles (\texttt{month, day}) : \textbf{one-hot encoding}
    \item Impl\'ementation : \texttt{Pipeline} + \texttt{ColumnTransformer} (scikit-learn)
  \end{itemize}

  \vspace{0.4em}
  Avantage : pr\'e-traitement inclus dans la validation et la recherche d'hyperparam\`etres (pas de fuite d'information).
\end{frame}

\begin{frame}{M\'ethodologie : split, validation et m\'etriques}
  \begin{itemize}
    \item Split train/test : 80\% / 20\%
    \item \textbf{Validation crois\'ee 10-fold} sur le train (robustesse)
    \item R\'egression : \RMSE/\MAE sur \(\LogArea\)
    \item Classification : Accuracy + Precision/Recall/F1 (split stratifi\'e)
  \end{itemize}
\end{frame}

\section{R\'egression}

\begin{frame}{Protocole exp\'erimental (r\'egression)}
  \begin{itemize}
    \item Baseline : \texttt{DummyRegressor} (m\'ediane)
    \item Mod\`eles : Linear Regression, SVR (RBF), KNN Regressor
    \item S\'election d'hyperparam\`etres : \texttt{GridSearchCV} avec \textbf{CV=10}
    \item Score optimis\'e : \texttt{neg\_root\_mean\_squared\_error}
  \end{itemize}
\end{frame}

\begin{frame}{R\'egression : GridSearch (SVR / KNN)}
  \begin{itemize}
    \item SVR (RBF) : recherche sur \(C\), \(\gamma\), \(\epsilon\)
    \item KNN : recherche sur \texttt{n\_neighbors}, \texttt{weights}, distance (\(p=1/2\))
    \item Objectif : am\'eliorer la performance vs baseline, sans sur-apprentissage
  \end{itemize}
\end{frame}

\begin{frame}{R\'esultats r\'egression (\RMSE/\MAE sur \(\LogArea\))}
  \textit{Note :} la cellule \textbf{Export Beamer} du notebook g\'en\`ere automatiquement ce tableau.

  \vspace{0.6em}
  \begin{table}
    \centering
    \small
    \IfFileExists{tables/regression_table.tex}{
      \begin{tabular}{lcc}
\toprule
Modèle & RMSE (test, log) & CV 10-fold RMSE (log) \\
\midrule
KNNRegressor(best) & 1.490 & 1.352 $\pm$ 0.185 \\
LinearRegression & 1.517 & 1.509 $\pm$ 0.340 \\
SVR(best) & 1.572 & 1.442 $\pm$ 0.244 \\
DummyRegressor(median) & 1.643 & 1.519 $\pm$ 0.279 \\
\bottomrule
\end{tabular}
    }{
      \begin{tabular}{lcc}
        \toprule
        Mod\`ele & RMSE (test, log) & CV 10-fold RMSE (log) \\
        \midrule
        DummyRegressor (median) & \textit{...} & \textit{...} \\
        LinearRegression & \textit{...} & \textit{...} \\
        SVR (best) & \textit{...} & \textit{...} \\
        KNNRegressor (best) & \textit{...} & \textit{...} \\
        \bottomrule
      \end{tabular}
    }
  \end{table}

  \vspace{0.2em}
  Message cl\'e : comparer \`a la baseline et juger la stabilit\'e via la CV.
\end{frame}

\begin{frame}{Discussion r\'egression (interpr\'etation)}
  \begin{itemize}
    \item Les mod\`eles non-lin\'eaires peuvent r\'eduire le RMSE (log), mais les gains restent souvent limit\'es
    \item La CV permet d'\'evaluer la stabilit\'e : forte variabilit\'e \(\Rightarrow\) risque de sur-apprentissage
    \item Les \textbf{\'ev\'enements extr\^emes} (grands feux) sont rares \(\Rightarrow\) pr\'ediction pr\'ecise difficile
  \end{itemize}
\end{frame}

\section{Classification}

\begin{frame}{Variante classification binaire}
  On transforme :
  \begin{itemize}
    \item Classe 1 : $\Area>0$ (un feu)
    \item Classe 0 : $\Area=0$ (aucun feu)
  \end{itemize}

  \vspace{0.6em}
  \begin{itemize}
    \item Split stratifi\'e train/test
    \item M\'etriques : Accuracy + Precision/Recall/F1
    \item CV : \textbf{Stratified 10-fold} sur le train
    \item Mod\`eles : Dummy, Logistic Regression, SVC (RBF), KNN
  \end{itemize}
\end{frame}

\begin{frame}{Classification : enjeu et d\'es\'equilibre}
  \begin{itemize}
    \item La cible $\Area>0$ peut \^etre d\'es\'equilibr\'ee (beaucoup de z\'eros)
    \item On privil\'egie F1 / Precision / Recall plut\^ot que l'Accuracy seule
    \item CV : \textbf{Stratified 10-fold} sur le train
  \end{itemize}
\end{frame}

\begin{frame}{R\'esultats classification (test + CV 10-fold)}
  \textit{Note :} la cellule \textbf{Export Beamer} g\'en\`ere aussi ce tableau.

  \vspace{0.2em}
  \begin{table}
    \centering
    \scriptsize
    \IfFileExists{tables/classification_table.tex}{
      \begin{tabular}{lcc}
\toprule
Modèle & F1 (test) & CV 10-fold F1 \\
\midrule
SVC(rbf) & 0.708 & 0.583 $\pm$ 0.052 \\
KNNClassifier & 0.689 & 0.618 $\pm$ 0.049 \\
DummyClassifier(most\_frequent) & 0.684 & 0.687 $\pm$ 0.009 \\
LogisticRegression & 0.593 & 0.494 $\pm$ 0.050 \\
\bottomrule
\end{tabular}
    }{
      \begin{tabular}{lcc}
        \toprule
        Mod\`ele & F1 (test) & CV 10-fold F1 \\
        \midrule
        DummyClassifier (most\_frequent) & \textit{...} & \textit{...} \\
        LogisticRegression & \textit{...} & \textit{...} \\
        SVC (rbf) & \textit{...} & \textit{...} \\
        KNNClassifier & \textit{...} & \textit{...} \\
        \bottomrule
      \end{tabular}
    }
  \end{table}

  \vspace{0.2em}
  \textbf{Option :} matrice de confusion export\'ee.
  
  \IfFileExists{figures/confusion_matrix_best.png}{
    \FFigure[height=0.30\textheight]{figures/confusion_matrix_best.png}
  }{}
\end{frame}

\begin{frame}{Matrice de confusion (meilleur mod\`ele)}
  \IfFileExists{figures/confusion_matrix_best.png}{
    \FFigure[height=0.68\textheight]{figures/confusion_matrix_best.png}
  }{\textit{(Figure non disponible : ex\'ecuter la cellule Export Beamer)} }

  \vspace{0.3em}
  Lecture : compromis faux positifs / faux n\'egatifs \(\Rightarrow\) choix de seuil selon l'objectif (Recall vs Precision).
\end{frame}

\begin{frame}{Discussion classification}
  \begin{itemize}
    \item Un bon F1 signifie un compromis correct entre d\'etection (Recall) et fausses alertes (Precision)
    \item La matrice de confusion permet de visualiser les erreurs syst\'ematiques
    \item Limite : la pr\'ediction \`a partir d'indices m\'et\'eo/FWI reste bruit\'ee (facteurs non observ\'es : intervention, topographie fine, etc.)
  \end{itemize}
\end{frame}

\section{Bonus \& Conclusion}

\begin{frame}{Bonus : s\'election de variables (top-4)}
  Approche (dans le notebook) : \textbf{Permutation Importance} sur un split interne du train.
  \begin{itemize}
    \item Objectif : garder seulement 4 variables pertinentes
    \item R\'e-entra\^iner un mod\`ele simple (ex: Linear Regression) avec ces variables
  \end{itemize}

  \vspace{0.4em}
  \IfFileExists{bonus_auto.tex}{\textbf{Top-4 variables (Permutation Importance) :} \texttt{DC}, \texttt{month}, \texttt{DMC}, \texttt{temp}\\
\textbf{RMSE (hectares) :} LinearRegression(raw)=107.767 ; LinearRegression(log→area)=109.799.}{\textbf{\`A compl\'eter :} liste des top-4 trouv\'ees.}
\end{frame}

\begin{frame}{Bonus : avec vs sans transformation log}
  Comparaison \`a interpr\'eter avec prudence :
  \begin{itemize}
    \item Sans log : l'erreur (en ha) est domin\'ee par les grands feux
    \item Avec log : apprentissage plus stable ; reconversion \(\widehat{\Area}=\exp(\hat{y})-1\) potentiellement biais\'ee
  \end{itemize}

  \vspace{0.4em}
  \textbf{\`A compl\'eter :} RMSE en hectares pour les deux variantes (cf. cellule Bonus 2).
\end{frame}

\begin{frame}{Conclusion}
  \begin{itemize}
    \item Pr\'edire $\Area$ (colonne \AreaCol) est difficile (beaucoup de z\'eros + extr\^emes)
    \item \(\LogArea\) am\'eliore la stabilit\'e pour la r\'egression
    \item Les mod\`eles non-lin\'eaires peuvent aider, mais la CV 10-fold est essentielle pour juger la robustesse
    \item En classification, l'Accuracy seule est insuffisante : regarder F1 + matrice de confusion
  \end{itemize}

  \vspace{0.6em}
  \textit{Limites / pistes :} features saisonni\`eres, ensembles, calibration/seuils (prioriser recall vs precision).
\end{frame}

\begin{frame}{R\'ef\'erences et reproductibilit\'e}
  \begin{itemize}
    \item Notebook d'analyse : \texttt{ForestFirePred.ipynb}
    \item Donn\'ees : \texttt{forestfires.csv}
    \item Code : scikit-learn (pipelines, GridSearchCV, cross\_validate)
  \end{itemize}

  \vspace{0.6em}
  Pour compiler : \texttt{pdflatex ForestFirePred\_presentation.tex}
\end{frame}

\end{document}
